%% 
%% Copyright 2007-2024 Elsevier Ltd
%% 
%% This file is part of the 'Elsarticle Bundle'.
%% ---------------------------------------------
%% 
%% It may be distributed under the conditions of the LaTeX Project Public
%% License, either version 1.3 of this license or (at your option) any
%% later version.  The latest version of this license is in
%%    http://www.latex-project.org/lppl.txt
%% and version 1.3 or later is part of all distributions of LaTeX
%% version 1999/12/01 or later.
%% 
%% The list of all files belonging to the 'Elsarticle Bundle' is
%% given in the file `manifest.txt'.
%% 
%% Template article for Elsevier's document class `elsarticle'
%% with numbered style bibliographic references
%% SP 2008/03/01
%% $Id: elsarticle-template-num.tex 249 2024-04-06 10:51:24Z rishi $
%%
%%\documentclass[preprint,12pt]{elsarticle}

%% Use the option review to obtain double line spacing
%% \documentclass[authoryear,preprint,review,12pt]{elsarticle}

%% Use the options 1p,twocolumn; 3p; 3p,twocolumn; 5p; or 5p,twocolumn
%% for a journal layout:
%% \documentclass[final,1p,times]{elsarticle}
%%\documentclass[final,1p,times,twocolumn]{elsarticle}
%% \documentclass[final,3p,times]{elsarticle}
\documentclass[final,3p,times,twocolumn]{elsarticle}
%% \documentclass[final,5p,times]{elsarticle}
%% \documentclass[final,5p,times,twocolumn]{elsarticle}

%% For including figures, graphicx.sty has been loaded in
%% elsarticle.cls. If you prefer to use the old commands
%% please give \usepackage{epsfig}

%% The amssymb package provides various useful mathematical symbols
\newcommand{\cmark}{\textcolor{green!60!black}{\checkmark}}
\newcommand{\xmark}{\textcolor{red!60!black}{\ding{55}}}
\usepackage{pifont}
\usepackage{threeparttable}
% \usepackage{multirow}
% \DeclareMathOperator*{\argmax}{arg\,max}
\usepackage{booktabs}
\usepackage{amssymb}
%% The amsmath package provides various useful equation environments.
\usepackage{amsmath}
\usepackage{graphicx}
\usepackage{multirow}%
\usepackage{amsmath,amssymb,amsfonts}%
\usepackage{amsthm}%
\usepackage{mathrsfs}%
\usepackage[title]{appendix}%
\usepackage{xcolor}%
\usepackage{textcomp}%
\usepackage{manyfoot}%
\usepackage{booktabs}%
\usepackage{url}
\usepackage[linesnumbered,ruled,vlined]{algorithm2e}
\usepackage{algorithmicx}%
\usepackage{algpseudocode}%
\usepackage{listings}%
\usepackage{float}
\usepackage{subcaption}
\usepackage[utf8]{inputenc}
\usepackage[T5]{fontenc}
\usepackage{appendix}
\usepackage[backref=page]{hyperref}
\hypersetup{
    colorlinks=true,
    linkcolor=orange,
    citecolor=blue,
    urlcolor=blue
}
%% The amsthm package provides extended theorem environments
%% \usepackage{amsthm}

%% The lineno packages adds line numbers. Start line numbering with
%% \begin{linenumbers}, end it with \end{linenumbers}. Or switch it on
%% for the whole article with \linenumbers.
%% \usepackage{lineno}

%\journal{Nuclear Physics B}

\begin{document}

\begin{frontmatter}

%% Title, authors and addresses

%% use the tnoteref command within \title for footnotes;
%% use the tnotetext command for theassociated footnote;
%% use the fnref command within \author or \affiliation for footnotes;
%% use the fntext command for theassociated footnote;
%% use the corref command within \author for corresponding author footnotes;
%% use the cortext command for theassociated footnote;
%% use the ead command for the email address,
%% and the form \ead[url] for the home page:
%% \title{Title\tnoteref{label1}}
%% \tnotetext[label1]{}
%% \author{Name\corref{cor1}\fnref{label2}}
%% \ead{email address}
%% \ead[url]{home page}
%% \fntext[label2]{}
%% \cortext[cor1]{}
%% \affiliation{organization={},
%%             addressline={},
%%             city={},
%%             postcode={},
%%             state={},
%%             country={}}
%% \fntext[label3]{}

% \title{Lightweight Data Augmentation Techniques for Imbalanced Learning}

% %% use optional labels to link authors explicitly to addresses:
% % \author[label1,label2]{}
% % \affiliation[label1]{organization={},
% %             addressline={},
% %             city={},
% %             postcode={},
% %             state={},
% %             country={}}

% % \affiliation[label2]{organization={},
% %             addressline={},
% %             city={},
% %             postcode={},
% %             state={},
% %             country={}}

% % =========================
% % Footnotes with symbols
% % =========================
% \fntext[dagger]{These authors contributed equally to this work (co-first authors).}
% \cortext[star]{Corresponding authors.}

% % =========================
% % Authors
% % =========================

% % ---- Co-first authors (†) ----
% \author[label1,fnref{dagger}]{Thai-Duy Dinh}
% \ead{dinhthaiduy.tk10@gmail.com}

% \author[label1,fnref{dagger}]{Minh-Thuan Nguyen}
% \ead{thuan2035@gmail.com}

% % ---- Middle authors ----
% \author[label1]{Thanh-Kim Le Duc}
% \ead{leducthanhkim2004@gmail.com}

% \author[label1]{Man-Dat Nguyen}
% %\ead{10421008@student.vgu.edu.vn}

% \author[label2]{Duc-Quang Vu}
% \ead{vdquang@ictu.edu.vn}

% \author[label3]{Ngoc-Hoang Tran}
% \ead{hoang2531992@gmail.com}

% \author[label4]{Anh-Cang Phan}
% \ead{cangpa@vlute.edu.vn}

% \author[label5]{Ngoc-Hoang Le}
% \ead{lengochoang252@gmail.com}

% \author[label6]{Phung-Anh Nguyen}
% \ead{alexnthhp@tmu.edu.tw}

% % ---- Co-corresponding authors (*) ----
% \author[label7]{Tan-Ha Mai\corref{cor}}
% \ead{maitanhaksdtvt6@gmail.com}

% \author[label1]{Lam-Son L\^{e}\corref{cor}}
% \ead{son.ll@vgu.edu.vn}

% % =========================
% % Affiliations
% % =========================
% \affiliation[label1]{organization={Faculty of Engineering, Vietnamese-German University},
%             city={Ho Chi Minh City},
%             country={Vietnam}}

% \affiliation[label2]{organization={Faculty of Information Technology,
% Thai Nguyen University of Information and Communication Technology},
%             city={Thai Nguyen},
%             country={Vietnam}}

% \affiliation[label3]{organization={Department of Software Engineering, FPT University},
%             city={Can Tho},
%             country={Vietnam}}

% \affiliation[label4]{organization={Vinh Long University of Technology Education},
%             city={Vinh Long},
%             country={Vietnam}}

% \affiliation[label5]{organization={Graduate Institute of Biomedical Materials and Tissue Engineering,
% College of Biomedical Engineering, Taipei Medical University},
%             city={Taipei},
%             country={Taiwan}}

% \affiliation[label6]{organization={Graduate Institute of Data Science,
% Taipei Medical University},
%             city={Taipei},
%             country={Taiwan}}

% \affiliation[label7]{organization={Department of Computer Science and Information Engineering,
% National Taiwan University},
%             city={Taipei},
%             country={Taiwan}}

% \cortext[cor]{Corresponding author}
\title{Lightweight Data Augmentation Techniques for Imbalanced Learning}

% =========================
% Footnotes with symbols
% =========================
\fntext[dagger]{These authors contributed equally to this work (co-first authors).}
\cortext[cor]{Corresponding authors.}

% =========================
% Authors
% =========================

% ---- Co-first authors (†) ----
\author[label1,fnref{dagger}]{Thai-Duy Dinh}
\ead{dinhthaiduy.tk10@gmail.com}

\author[label1,fnref{dagger}]{Minh-Thuan Nguyen}
\ead{thuan2035@gmail.com}

% ---- Middle authors ----
\author[label1]{Thanh-Kim Le Duc}
\ead{leducthanhkim2004@gmail.com}

%\author[label1]{Man-Dat Nguyen}

\author[label2]{Duc-Quang Vu}
\ead{vdquang@ictu.edu.vn}

\author[label3]{Ngoc-Hoang Tran}
\ead{hoang2531992@gmail.com}

\author[label4]{Anh-Cang Phan}
\ead{cangpa@vlute.edu.vn}

\author[label5]{Ngoc-Hoang Le}
\ead{lengochoang252@gmail.com}

\author[label6]{Phung-Anh Nguyen}
\ead{alexnthhp@tmu.edu.tw}

% ---- Co-corresponding authors (*) ----
\author[label7]{Tan-Ha Mai\corref{cor}}
\ead{maitanhaksdtvt6@gmail.com}

\author[label1]{Lam-Son L\^{e}\corref{cor}}
\ead{son.ll@vgu.edu.vn}

% =========================
% Affiliations
% =========================
\affiliation[label1]{organization={Faculty of Engineering, Vietnamese-German University},
            city={Ho Chi Minh City},
            country={Vietnam}}

\affiliation[label2]{organization={Faculty of Information Technology,
Thai Nguyen University of Information and Communication Technology},
            city={Thai Nguyen},
            country={Vietnam}}

\affiliation[label3]{organization={Department of Software Engineering, FPT University},
            city={Can Tho},
            country={Vietnam}}

\affiliation[label4]{organization={Vinh Long University of Technology Education},
            city={Vinh Long},
            country={Vietnam}}

\affiliation[label5]{organization={Graduate Institute of Biomedical Materials and Tissue Engineering, Taipei Medical University},
            city={Taipei},
            country={Taiwan}}

\affiliation[label6]{organization={Graduate Institute of Data Science, Taipei Medical University},
            city={Taipei},
            country={Taiwan}}

\affiliation[label7]{organization={Department of Computer Science and Information Engineering,
National Taiwan University},
            city={Taipei},
            country={Taiwan}}

%% Abstract
\begin{abstract}
    Imbalanced learning, when some classes have far fewer examples than others, creates real challenges in areas like medical diagnosis, fraud detection, and rare-event forecasting. Advanced generative methods such as GANs, diffusion models, and DeepSMOTE, can effectively augment minority-class examples but typically demand substantial compute resources and complex pipelines, which limits their practicality in environments with constrained resources. In this work, we focus instead on lightweight data-augmentation techniques that balance speed and accuracy. By combining simple augmentations (Mixup, Cutout) with careful resampling and class-weighting, we achieve significant gains in minority-class performance without extra computational cost. Evaluated on imbalanced versions of SVHN-10, CIFAR-10, CINIC-10, CIFAR-100 and TinyImageNet-200 our approach consistently improves both minority-class accuracy and overall generalization, while keeping training fast and resource-friendly. These results offer a practical solution for handling class imbalance in constrained environments. Our implementation is open-sourced and available on GitHub\footnote{\url{https://github.com/duydinhthai27/lightweight-ILWH}} to foster future research in this direction.
\end{abstract}



%%Graphical abstract
% \begin{graphicalabstract}
% %\includegraphics{grabs}
% \end{graphicalabstract}

%%Research highlights
% \begin{highlights}
% \item Research highlight 1
% \item Research highlight 2
% \end{highlights}

%% Keywords
\begin{keyword}
Deep Learning, Imbalanced Classification, Lightweight Data Augmentation, Imbalanced Lightweight Hybrid Approach.
\end{keyword}

\end{frontmatter}

%% Add \usepackage{lineno} before \begin{document} and uncomment 
%% following line to enable line numbers
%% \linenumbers

%% main text
%%

%% Use \section commands to start a section
\input{introduction}
\input{related_work}
\input{why-lightweight}
\input{Lightweight_imbalanced}
\input{ablation-study}
\input{experiment}
\input{conclusion}
\input{limitation}

% \clearpage
\bibliography{reference.bib}
\bibliographystyle{elsarticle-num}

\end{document}

\endinput
%%
%% End of file `elsarticle-template-num.tex'.
